% Generated draft: requires mathematical and source review.
\documentclass[11pt]{article}
\usepackage[margin=25mm]{geometry}
\usepackage{fontspec}
\setmainfont{Latin Modern Roman}
\usepackage{amsmath,amssymb,mathtools,unicode-math}
\setmathfont{Latin Modern Math}
\usepackage{enumitem,needspace,xurl,hyperref}
\setlength{\parindent}{0pt}
\setlength{\parskip}{5pt}
\setlength{\emergencystretch}{3em}
\newcommand{\N}{\mathbb{N}}
\newcommand{\Z}{\mathbb{Z}}
\newcommand{\Q}{\mathbb{Q}}
\newcommand{\R}{\mathbb{R}}
\newcommand{\C}{\mathbb{C}}
\newcommand{\F}{\mathbb{F}}
\newcommand{\A}{\mathbb{A}}
\newcommand{\PP}{\mathbb{P}}
\newcommand{\E}{\mathbb{E}}
\newcommand{\eps}{\varepsilon}
\newcommand{\dd}{\,\mathrm{d}}
\begin{document}
\section*{UnsolvedMath ID: 30001260}
\textbf{Source problem code:} catalogue entry
\subsection{Definitions and mathematical statement}
Let \(S_5\) be the group of permutations of \(\{1,2,3,4,5\}\). For an integer \(n\geq 0\), write \(S^n=\{x\in\R^{n+1}:\|x\|=1\}\), with its standard smooth structure. A smooth action is a homomorphism \(\rho:S_5\to\operatorname{Diff}(S^n)\). Its isotropy subgroup at \(x\in S^n\) is
\[
(S_5)_x=\{g\in S_5:\rho(g)(x)=x\}.
\]
\textbf{Supplied question.} Does \(S_5\) admit a smooth action on a sphere with rank-one isotropy?

The supplied record does not define ``rank-one isotropy.'' One possible rank convention for a finite group \(H\) is its elementary abelian rank:
\[
r_p(H)=\max\{k\in\Z_{\geq0}:(\Z/p\Z)^k\text{ embeds in }H\},\qquad
r(H)=\max_{p\text{ prime}}r_p(H).
\]
Under this convention, the question would ask for a dimension \(n\) and a smooth action satisfying either
\[
\forall x\in S^n,\quad r((S_5)_x)=1,
\]
or, if ``rank-one isotropy'' means isotropy of rank at most one,
\[
\forall x\in S^n,\quad r((S_5)_x)\leq1.
\]
These are alternative conditional formulations: the second permits trivial stabilizers, whereas the first excludes them. The supplied evidence does not determine the rank convention or which condition is intended. It also does not specify a dimension restriction, whether the action must be effective, or any convention concerning the smooth structure of the sphere. The standard-sphere notation above makes the conditional formulations explicit; those unresolved qualifications require checking against the report before selecting a definitive statement. No existence or nonexistence assertion is made.
\subsection{Short English statement}
The question asks whether the permutation group \(S_5\) can act smoothly on some sphere with rank-one stabilizers. The meaning of rank-one isotropy, including whether trivial stabilizers are allowed, is not specified in the supplied record and requires source clarification.
\subsection{Sources}
\begin{itemize}
\item Local UnsolvedMath registry, permanent ID \(30001260\), source problem code \texttt{catalogue entry}, title: \emph{Rank-One-Isotropy Actions of \(S_5\) on Spheres}. The record belongs to the Oberwolfach Reports Open Problems collection and gives proposed year 2009; no individual proposer is supplied.
\item \emph{Manifold Perspectives} (2009), Oberwolfach Reports, pp.\ 1487--1546, DOI \texttt{10.4171/owr/2009/27}: \url{https://doi.org/10.4171/owr/2009/27}. This is the supplied source citation and the supplied literature-source link, labeled ``original Oberwolfach report.'' The background describes these reports as documenting mathematical workshops and their open-problem discussions.
\item Supplied literature assessment, checked \texttt{2026-08-21}: the record labels the problem open, states that the report explicitly calls this question open, and reports that no later resolving construction was identified. This is supplied dated curation, not an independently verified assessment. The supplied triage expressly requires independent verification of time-sensitive status and cited claims.
\item The registry labels the statement \texttt{corrected\_verified}; its accompanying verification note says that mathematical wording was retained while notation and sentence punctuation were normalized for valid LaTeX rendering. This supplied label does not resolve the formulation issues identified here or establish mathematical review.
\end{itemize}
\end{document}
